
\begin{frame}{XPath}
XPath est un langage \alert{d'expression de chemins}.

\vfill

Il permet l'accès à des 
 nœuds d’un document XML en décrivant les chemins menant à ces nœuds.

\vfill

XPath est un outil de base dans la manipulation de documents XML. Utilisé dans 
\begin{description}
\item[XQuery] XPath fait partie du langage d’interrogation XQuery,
\item[XSLT] pour la sélection de la partie des données à transformer,
\item[XML Schema] pour les expressions de contraintes d’unicité,
%\item[XSL] pour la sélections des parties des données à formater,
\item[XPointer] pour l'identification de fragments,
\item[XLink] pour ancrer les liens hypertextes,
\item[...]
\end{description}

Ce qui est montré dans ce qui suit: XPath 1.0

\end{frame}


\section{Un exemple commenté}

\begin{frame}{Une intuition}

Permet de ``pointer'' (localiser) des nœuds dans un document XML.

\begin{exampleblock}{Intuition}
Par \alert{navigation} : en se déplaçant dans l'arbre de "étape par étape" à partir de la racine, on va chercher à 
   atteindre certaines de ses parties.
\end{exampleblock}

\vfill

\begin{block}{Résultat d'une expression XPath}
Le résultat de l'évaluation d'une expression XPath est un \alert{ensemble de nœuds} sélectionnés dans le
document en entrée.
\end{block}

\vfill

\alert{Attention}.  Un nœud peut être racine d'un sous-arbre; dans ce cas
c'est l'ensemble du sous-arbre qui fait partie du résultat.
 
\end{frame}

\begin{frame}[allowframebreaks]{Commen\c{c}ons par un exemple: \texttt{/film/titre/text()}}

Vous vous souvenez sans doute de notre document XML représentant le film Gravity, dans
sa forme arborescente. Notez bien les types DOM des nœuds

\figSlideWithSize{GravityDOM}{9cm}

\vfill


\begin{remark}
Dans ce qui suit le nœud courant est marqué par des \alert{bords rouges en pointillés},
les nœuds sélectionnés par des \alert{bords rouges pleins} et les chemins par
des flèches. 
\end{remark}

\framebreak

Notre expression commence par un '/': c'est une \alert{expression absolue} qui prend
pour nœud courant initial la \alert{racine du document}. 

\vfill

\figSlideWithSize{GravityDOM-xpath1}{9cm}

\vfill
\begin{block}{La navigation}
On navigue vers le bas, vers les enfants du nœud courant.
\end{block}

\framebreak

Une expression XPath est constituée \alert{d'étapes}. La première étape
est \texttt{film}: parmi les enfants du nœud courant, on prend
celui dont le \alert{nom} est \texttt{film}.
 
\vfill


\figSlideWithSize{GravityDOM-xpath2}{9cm}

\vfill

\vfill
\begin{block}{Les tests de nœud}
La sélection d'un nœud par son nom est un exemple de \alert{test de nœud} (\textit{node test}). 
Avec XPath on peut appliquer des filtres sur la \alert{structure} avec les \alert{tests de nœuds} et
sur le \alert{contenu} avec les \alert{prédicats}.
\end{block}

\framebreak

On passe à l'étape suivante: \texttt{titre}. Maintenant, le \alert{nœud courant} est
\texttt{film}, résultat de l'étape précédente. On regarde les enfants et on sélectionne
celui dont le nom est \texttt{titre}

\figSlideWithSize{GravityDOM-xpath3}{9cm}


\vfill
\begin{block}{Etape et nœud courant}
Le nœud courant à l'étape $n$ est un nœud sélectionné à l'étape $n-1$. On répète l'étape $n$
autant de fois qu'il y a de nœuds sélectionnés.
\end{block}

\framebreak

On arrive au bout\,! La dernière étape est \texttt{text()}. C'est un test de nœud sur le
\alert{type} du nœud, et non sur le nom.

\vfill

On prend donc, parmi les enfants du nœud courant \texttt{titre}, ceux de type \textbf{Text}.

\vfill

\figSlideWithSize{GravityDOM-xpath4}{9cm}

\vfill

C'est fini: l'expression a donné  chemin 
qui va de la racine du document vers le nœud-texte \texttt{Gravity}.

\end{frame}

\section{Formalisme XPath, en bref}

\subsection{Etapes, expressions, interprétation}

\begin{frame}{Qu'est-ce qu'une \alert{étape} XPath}

Une \alert{étape}  est une expression de la forme 
\texttt{\textcolor{indiaRed}{axis}::\textcolor{DarkOliveGreen}{node-test}\textcolor{cyan}{[predicate]}
}\\


\begin{description}
\item[\texttt{\textcolor{indiaRed}{axis}}] est la direction de la navigation\\

   Par défaut, on descend vers les enfants, mais beaucoup d'autres axes sont possibles.
   
\item[\texttt{\textcolor{DarkOliveGreen}{node-test}}] est  un test de structure\\

   Il porte sur le type (DOM) du nœud ou sur le nom du nœud (éléments et attributs seulement)
   
\item[\texttt{\textcolor{cyan}{[predicate]}}] (optionnel) est un test sur le contenu\\

    On prend les valeurs de nœuds-texte ou des attributs, on les compare à des constantes ou entre eux:
    un peu à la SQL.
    
\end{description}

\medskip

Dans \texttt{child::realisateur[@nom='Cuaron']}, 
\texttt{child} est l'axe, \texttt{library} le test,  le prédicat est entre crochets. 

\end{frame}


\begin{frame}{Interprétation d'une expression}

\begin{block}{Principe général}
Une expression XPath est une suite d'étapes (\textit{steps}), chaque étape comprenant éventuellement des tests 
sur les nœuds atteints pour en éliminer certains.
\end{block}

\medskip
A chaque étape\,:
\begin{enumerate}
  \item on considère un \alert{nœud courant}, pris comme point de départ\,;
   \item on ``désigne'' à partir du nœud courant un ensemble d'autres nœuds\,;
   \item on applique à chacun les tests de nœud\,;
   \item on applique les prédicats (optionnels)\,;
   \item enfin on prend chacun des nœuds restant comme nœud courant pour l'étape suivante.
\end{enumerate}
\end{frame}


\subsection{Les axes XPath}

\begin{frame}{Les axes dans XPath}

À partir d'un nœud de l'arbre (nœud courant), les axes permettent 
de se déplacer dans tous les sens.   

\vfill

On va se contenter du minimum (qui est aussi le plus utile!)

\begin{itemize}
  \item L'axe \texttt{child}: les enfants du nœud courant.
  \item \texttt{descendant}: tous les descendants du nœud courant.
  \item \texttt{parent}: l'ascendant direct du nœud courant. 
  \item \texttt{attribute}: les attributs du nœud courant.
  \item \texttt{self}: le nœud courant lui-même.
\end{itemize}

\begin{remark}
Si vous êtes amenés à utiliser XPath intensivement, il faudra aller plus loin:
reportez-vous au livre \url{http://webdam.inria.fr/jorge} pour une présentation complète.
\end{remark}
\end{frame}


\begin{frame}{L'axe \texttt{child}}

Pas de surprise: à chaque \alert{étape}, on se dirige vers les enfants du nœud courant.

\vfill

La figure montre les nœuds désignés par l'expression \texttt{child::node()}, pour le
nœud courant \texttt{film}.

\figSlideWithSize{GravityDOM-child}{9cm}

\vfill

\begin{remark}
Note: \texttt{node()} est le test "neutre": il renvoie vrai pour tous les nœuds (ici, sans
tenir compte du \alert{nom} du nœud).
\end{remark}
\end{frame}


\begin{frame}{L'axe \texttt{descendant}}

On parcourt les enfants, puis les enfants des enfants, jusqu'aux feuilles.

\vfill

La figure montre les nœuds désignés par l'expression \texttt{descendant::node()}, pour le
nœud courant \texttt{film}.

\figSlideWithSize{GravityDOM-descendant}{9cm}

\begin{remark}
Si, quand même une petite surprise: les \alert{attributs} ne sont pas concernés. 
\end{remark}
\end{frame}


\begin{frame}{L'axe \texttt{attribute}}

Les attributs sont spéciaux: on ne peut y accéder qu'en utilisant explicitement l'axe \texttt{@}
(ou \texttt{attribute}).

\vfill

La figure montre les nœuds désignés par l'expression \texttt{@*}, pour le
nœud courant \texttt{réalisateur}.

\figSlideWithSize{GravityDOM-attribute}{9cm}

\end{frame}

\section{Tests et prédicats}


\subsection{Test de nœud}

\begin{frame}{Les tests de nœuds}

Un axe désigne un ensemble de nœuds (par exemple, les enfants du nœud courant).

\vfill

On les \alert{filtre} avec un test (\texttt{node-test}). 

\vfill

On peut tester sur le type (DOM) du nœud :
\begin{description}
\item[\texttt{node()}] sélectionne les nœuds  quel que soit leur type
\item[\texttt{text()}] sélectionne les nœuds de type \textbf{Text}.
\item[\texttt{element()}] sélectionne les nœuds de type \textbf{Element}.
\item et quelques autres...
\end{description}

\medskip

Ou bien tester sur le nom du nœud (seulement pour éléments et attributs).

\begin{description}
\item[\texttt{leNom}] sélectionne les éléments ou attributs de nom \texttt{leNom}
\item[\texttt{*}] sélectionne les nœuds qui ont un nom, quel que soit le nom.
\end{description}

\end{frame}

\begin{frame}{Exemple: sélection des nœuds texte}

On veut tous les nœuds de type \textbf{Text}.

\vfill

Expression: \texttt{/film/*/text()}.

\figSlideWithSize{GravityDOM-text}{9cm}

\vfill

Quels sont les tests dans cette expression? Quels sont les axes?

\end{frame}




\subsection{Les prédicats}

\begin{frame}{Les prédicats}

Alors que les \texttt{node-test} portent sur la \alert{structure} du document,
les \alert{prédicats} portent sur le contenu.

\vfill

Quelques exemples, mieux que les discours.


\begin{itemize}
\item Le titre du film réalisé par M. Cuaron
 
	\texttt{/film[realisateur/@nom='Cuaron']/titre}
	
\item Les films de SF: \texttt{/film[genre='SF']/titre}

\item Les personnes dont on connaît la date de naissance:

    \texttt{/personne[@date\_naissance]}
    
 \item Les films dont le genre est soit SF, soit comédie:
 
   \texttt{/film[genre='SF' or genre='Comedy']/titre}
\end{itemize}

\begin{block}{Le principe}
Un nœud $n$ est sur un chemin. Pour savoir si on le garde ou pas: tests
sur des \alert{valeurs} ramenées par des expressions \alert{relatives} à $n$.  
\end{block}
\end{frame}


\begin{frame}{Un exemple détaillé}

Films réalisés par M. Cuaron: \texttt{/film[realisateur/@nom='Cuaron']/titre}

\vfill

\figSlideWithSize{GravityDOM-predicate}{9cm}

\vfill

\begin{block}{Ce qu'il vaut mieux cacher...}
Que se passe-t-il quand l'expression d'un prédicat désigne \alert{plusieurs} nœuds
que l'on compare à \alert{une} valeur? Passons pudiquement...
\end{block}

\end{frame}

\begin{frame}{Abréviations}

Certaines expressions très courantes ont une abréviation.

\vfill


\textcolor{red}{.} désigne le nœud courant (abrégé de \texttt{self::node()})

\vfill


\textcolor{red}{..} désigne le  parent du nœud courant (abrégé de \texttt{parent::node()})

\vfill

\textcolor{red}{\texttt{//}} est une expression abrégée désignant tous les descendants du 
  nœud courant, y compris le nœud courant lui-même.
  
  \begin{itemize}
     \item \textcolor{red}{\texttt{//a}} désigne tous les nœuds du document dont le nom est \texttt{a}.
     \item \textcolor{red}{\texttt{.//b}} désigne tous les descendants du nœud courant (ou lui-même)
          dont le nom est \texttt{b}.
  \end{itemize}


\end{frame}

\begin{frame}[allowframebreaks]{Opérations et fonctions}

XPath permet aussi d'effectuer des opérations et fonctions. Mini sélection:

\begin{itemize}
\item \texttt{count(nset)} nombre d'éléments dans l'ensemble de nœuds. 
\item \texttt{name(nœud)} nom du nœud.
\item \texttt{concat(ch1, ..., chn)}, concaténation de textes
\item \texttt{contains(ch1,ch2)}, vrai si \texttt{ch1} contient \texttt{ch2}.
\end{itemize}


{\bf Fonctions sur les booléens}

\begin{itemize}
\item \texttt{not(bool)}
\end{itemize}

\begin{remark}
La version 2 de XPath est un sous-ensemble de XQuery: toutes les fonctions XQuery
sont disponibles, il y en a beaucoup et c'est extensible.
\end{remark}
\end{frame}



\begin{frame}{Bilan}

Nous en savons assez pour faire beaucoup de choses intéressantes.

\begin{itemize}
  \item Tous les nœuds texte du document: \texttt{//text()}
  \item Les nœuds texte dont le nom du parent est \texttt{titre}:
       \texttt{//text()[parent::titre]}
         \item Les nœuds texte dont le nom du parent est \texttt{titre}
         et qui contiennent la sous-chaîne \texttt{'ty'}.
         
    \texttt{//text()[parent::titre and contains(., 'ty')]}
       
    \item Les nœuds qui ont plus de deux enfants.
    
      \texttt{//*[count(*) >= 2]} 
       
\end{itemize}

\begin{block}{Pour conclure}
Deux remarques
\begin{itemize}
  \item XPath sélectionne des nœuds existant; on ne peut pas construire de nouveaux documents,
      de nouvelles valeurs.
   \item Typique d'une approche où on peut interroger à la fois la structure et le contenu.   
\end{itemize}
\end{block}

\end{frame}